\begin{table}[t]
\caption{Recuperación del rendimiento y brecha respecto del entrenamiento limpio por familia de contramedida, en puntos porcentuales de exactitud}\label{tab:familias}
\footnotesize
\begin{tabularx}{\textwidth}{@{}L{3.0cm}*{6}{Z{1.0}}@{}}
\toprule
Familia de contramedida & Estudios & Comparaciones & Recuperación & Rango intercuartílico & Brecha al entrenamiento limpio & Cercanía al entrenamiento limpio (\%) \\
\midrule
Purificación generativa & 3 & 267 & 66,8 & -- & 7,9 & 32 \\
Cambio de paradigma de aprendizaje & 1 & 79 & 45,7 & -- & 14,4 & 19 \\
Aumentos adversarios & 2 & 31 & 38,4 & -- & 33,9 & 0 \\
Proyección y sondas lineales & 2 & 38 & 33,5 & -- & 35,9 & 4 \\
Entrenamiento adversario & 10 & 411 & 25,9 & 20,6--39,9 & 39,1 & 3 \\
Transformación y compresión de entrada & 12 & 1248 & 25,0 & 8,3--32,6 & 23,9 & 15 \\
Otras & 6 & 354 & 25,8 & 8,7--29,7 & 41,3 & 16 \\
Aumentos de datos & 8 & 642 & 4,7 & 0,5--15,5 & 66,4 & 4 \\
\botrule
\end{tabularx}
\footnotetext{Recuperación y brecha: mediana de las medianas por estudio. El rango intercuartílico se informa con cuatro o más estudios. Cercanía: porcentaje medio, por estudio, de comparaciones que quedan a 5 puntos o menos del entrenamiento limpio.}
\end{table}
